\section{Discussion}
\label{sec:discussion}

\subsection{How a mechanism can be active and uninformative at once}

The two findings that look contradictory---displacement of a third of the
population spread, and complete insensitivity to permutation---are reconciled
by \eqref{eq:leverage}. The displacement is governed by $\chi_i$, which depends
only on how \emph{concentrated} the weights are, not on where the concentration
falls. Betweenness on a sparse small-world graph is highly dispersed and
exactly zero at many nodes, so $\chi_i$ is large; the guard $\varepsilon$ then
makes the operator behave as hard masking, excluding a subset of neighbours
from each aggregate. Permuting the scores changes \emph{which} neighbours are
excluded but not \emph{how many}, leaving $\chi_i$ statistically unchanged.

What this says about the mechanism is specific: it perturbs the neighbourhood
aggregate strongly, and the perturbation is effectively a structured random
subsampling of neighbours. Subsampling neighbours has an effect---it is the
(P1) contribution we measure at $3$ of $28$ functions---but the identity of the
retained neighbours does not, because betweenness on a graph whose edges bear
no relation to the objective landscape cannot be related to the landscape
either. The graph is a small-world prior fixed before the first evaluation;
nothing in the search updates it toward the problem.

This suggests where a structural mechanism could genuinely carry information,
and it is not in the weighting but in the graph. A topology derived from the
landscape---edges reflecting estimated barrier heights or basin membership
rather than a Watts--Strogatz prior---would make betweenness a statement about
the search space rather than about an arbitrary lattice. Local optima
networks~\cite{ochoa2008lon} build exactly such graphs, but offline and for
characterisation; the wider landscape-analysis literature~\cite{malan2021survey}
is likewise diagnostic rather than generative. Whether an online estimate is accurate enough to drive
search is open, and the controls developed here are the instrument that would
settle it: the prediction is that on a landscape-derived graph, unlike the one
tested, permutation would hurt.

\subsection{A practical consequence}

Betweenness centrality is roughly $30\%$ of per-iteration runtime in the
reference implementation, and computing it costs $O(N\!\cdot\!M)$. Our results
say it can be replaced by any heterogeneous weighting---a fixed random draw
suffices---with no measurable loss, and by uniform weighting at a cost of about
$0.17$ rank units on the CEC2017 suite. For practitioners the recommendation is
concrete: if a centrality-weighted method is used, compute the cheapest
centrality available, or none.

\subsection{Density is the component that survives, and it resists prediction}

Of the structural choices we examined, neighbourhood size is the one with a
large and reproducible effect: consequential in $30$ of $49$ configurations,
with rank spreads up to $3.9$ of $5$. It is also the one the literature has
studied longest, from Kennedy and Mendes~\cite{kennedy2002population} through
the cellular-EA ratio~\cite{alba2005exploration}. Our contribution here is
negative and specific: the optimum is not predictable in advance from a
landscape ruggedness descriptor. The correlation is moderate
($\rho_{s}=-0.61$), carried by a few functions, absent at $D=10$, and
contradicted outright by the penalised functions.

That negative result is what motivated an adaptive scheme rather than a
predictive one, and it is worth stating why the distinction matters. A
predictor must be right about the landscape; an adaptive rule need only be
right about which of a few alternatives is currently doing better. The second
is a far weaker requirement, and our results suggest it is the achievable one:
the predictive route fails at $D=10$ and on the penalised functions, while the
adaptive route beats its own randomisation on $24$ of $28$ functions.

\subsection{On applying the instrument to our own proposal}

We built \textsc{Ads} and subjected it to the same policy control we would ask
of anyone else. It passes: informed selection beats uniform switching over the
same densities by $\bar{\delta}=-0.101$, $p=1.5\times10^{-5}$, on $24$ of $28$
functions, so the benefit is not merely from varying density, which both arms
share.

Reaching that conclusion required correcting our own first reading. Under
per-function Holm correction \textsc{Ads} wins on $3$ of $28$, which we
initially took for a weak effect; correction across $28$ comparisons is
conservative, and a benefit spread thinly across a suite fails it while being
obvious in aggregate. The diagnostic is the direction count, visible only when
the suite is the unit of analysis. Section~\ref{sec:target2:views} gives the
mirror case.

The contrasts together also say something about the optimizer under study. Its
two weighting mechanisms are null and weak at suite level, while the structural
choice that is not a weighting---how many neighbours each agent sees---carries
a clear effect. The information here lives in the coarse topology, not in the
scores computed on it.

\subsection{Implications for evaluation practice}

The confound we identify is structural, not incidental. Whenever a component
both alters the operator and is claimed to carry information, ablation reports
the sum of the two effects and the claim is attributed the whole. The remedy is
cheap: a permutation of the values the component consumes, or a randomisation
of the policy it implements, costs one additional arm and no additional budget
per arm. We would encourage reporting it wherever a structural or adaptive
claim is made, in the same way that effect sizes and non-parametric tests have
become expected.

Two caveats belong with that recommendation. First, a control is only
interpretable if the mechanism is active, which is why we report $\rho$
alongside every comparison; on a vertex-transitive topology the comparison is
vacuous by \eqref{eq:leverage}. Second, an equivalence claim needs an
equivalence test, and adequate power for it is more demanding than it looks:
equivalence at margin $m$ requires the point estimate \emph{and} its interval
to fit inside the margin, so halving the interval width does not suffice when
the estimate itself sits near the boundary. Our own per-function claims stop
at the small-effect margin for that reason, and we say so rather than resting
on non-significance.
